\section*{Introducción}
\justifying

La electrónica de potencia es una disciplina que permite transformar y controlar la energía eléctrica para
una amplia variedad de aplicaciones industriales, comerciales y residenciales. Aprender sobre estos
sistemas puede resultar complejo debido a la falta de acceso a equipos adecuados
y a la necesidad de una sólida comprensión teórica respaldada por la experimentación práctica.

El proyecto consiste en el diseño y construcción de un sistema didáctico modular con el objetivo de
facilitar la enseñanza y el aprendizaje de la electrónica de potencia, sin la necesidad de equipos
especializados y costosos. Este sistema permitirá a estudiantes o entusiastas experimentar con 
diversos módulos, tales como rectificadores, inversores, convertidores DC-DC, entre otros.

Mediante la observación de señales de entrada y salida y el ajuste de parámetros en tiempo real,
el sistema permitirá a los usuarios aprender y experimentar con diversos conceptos de
electrónica de potencia.

Además, contará con un manual de usuario que explique el funcionamiento teórico de cada módulo y
una guía práctica que oriente sobre los parámetros a modificar y los efectos a analizar. Con 
este enfoque, se espera que los usuarios desarrollen una comprensión profunda y tangible de los 
principios de la electrónica de potencia.

El diseño modular del sistema permitirá que, en el futuro, se puedan incorporar
expansiones o módulos personalizados según necesidades del usuario o avances tecnológicos.
De este modo, el proyecto queda abierto a mejoras continuas y a la adaptación a nuevas áreas
de estudio de la electrónica.